\documentclass[12pt]{article}
\usepackage[utf8]{inputenc}
\usepackage[T1]{fontenc}
\usepackage{hyperref}
\title{Inland-Buhi-Daraga Bikol Language Resources for K-12 Curriculum Development}
\author{Compiled by Claude Sonnet 5.5 (claude-sonnet-5-5) for Andr\'as Kornai}
\date{}
\begin{document}
\maketitle

\section{Introduction}
The code bhk (Glottocode alba1269) is an older label for Albay Bikol, a group of related Bikol varieties of the Philippines; translatewiki marks bhk as deprecated and split into distinct languages (including Buhi'non, Libon, West Albay and Miraya Bikol). Sources are not fully consistent on boundaries and codes. The macro-area given in the task (Papunesia) fits the Philippines; the family is Austronesian.

\section{Primary Educational Resources}
No Bikol-specific primers or textbooks were identified in the search results. The Philippines' mother-tongue-based education policy is not documented in the results found, so it is not cited here.

\section{Linguistic and Reference Materials}
\begin{itemize}
\item Wikipedia, Albay Bikol language: \url{https://en.wikipedia.org/wiki/Albay_Bikol_language}
\item Wikipedia, Bikol languages: \url{https://en.wikipedia.org/wiki/Bikol_languages}
\item Grambank entry for alba1269: \url{https://grambank.clld.org/languages/alba1269}
\item translatewiki, Bikol language list: \url{https://translatewiki.net/wiki/Languages_by_language_family/Bikol}
\end{itemize}

\section{Religious and Community Resources}
None found in the searches.

\section{Digital and Audio Resources}
None found in the searches.

\section{Curriculum Development Context}
Because the label is a deprecated umbrella, curriculum work should target a specific variety (e.g. Buhinon, Libon, Miraya, West Albay) and look for materials under those names; this search did not find them.

\end{document}